\documentclass[12pt]{article}

\usepackage[a4paper,margin=1in]{geometry}
\usepackage{amsmath,amssymb}
\usepackage{newtxtext,newtxmath}
\usepackage{fancyhdr}
\usepackage{listings}
\usepackage{graphicx}
\usepackage{float}
\usepackage{hyperref}
\usepackage{booktabs}
\usepackage{xcolor}

\setlength{\parindent}{0pt}
\setlength{\parskip}{5pt}

% Keep entire document at 12 pt
\AtBeginDocument{
    \fontsize{12pt}{14pt}\selectfont
}

% ------------------------------------------------
% Header Configuration
% ------------------------------------------------
\pagestyle{fancy}
\fancyhf{}
\fancyhead[L]{Name: \underline{Shaik Jnaved}}
\fancyhead[R]{Enrollment No.: \underline{P26DS001}}
\renewcommand{\headrulewidth}{0.4pt}
\setlength{\headheight}{15pt}

% ------------------------------------------------
% Code Formatting Configuration
% ------------------------------------------------
\lstset{
    basicstyle=\ttfamily\fontsize{8.5pt}{10.5pt}\selectfont,
    breaklines=true,
    frame=single,
    columns=fullflexible,
    keepspaces=true,
    showstringspaces=false,
    keywordstyle=\color{blue}\bfseries,
    commentstyle=\color{gray}\itshape,
    stringstyle=\color{teal}
}

\begin{document}

% ==============================================================================
% TITLE BANNER
% ==============================================================================
\begin{center}
\textbf{\Large SARDAR VALLABHBHAI NATIONAL INSTITUTE OF TECHNOLOGY, SURAT}\\[4pt]
\textbf{\large Department of Computer Science \& Engineering}\\[3pt]
\textbf{\large CS103 / CSDS103: Design and Analysis of Algorithms}\\[3pt]
\textbf{\large M.Tech. (DS) I Year -- Semester I}\\[4pt]
\textbf{\Large Lab Assignment 9: Advanced Graph Algorithms, Greedy Strategies, and Network Flows}
\end{center}

\vspace{0.2cm}
\hrule
\vspace{0.4cm}

\textbf{Overview \& Scope}:\\
This laboratory report presents rigorous mathematical formulation, design, complexity analysis, and empirical verification for six canonical algorithmic problems across graph theory, computational geometry, network flows, and greedy optimization:
\begin{enumerate}
    \item \textbf{Assignment 1}: Optimal Cryptographic Software License Purchase Schedule under Exponential Inflation via Greedy Ordering.
    \item \textbf{Assignment 2}: Linear-Time Directed Eulerian Circuit Construction using Hierholzer's Algorithm.
    \item \textbf{Assignment 3}: 3D Pick-up Sticks Geometric Relation Routine and Topological Sort Scheduling.
    \item \textbf{Assignment 4}: Linear-Time Two-Colorability (Bipartiteness) Testing and Conflict-Free Graph Partitioning.
    \item \textbf{Assignment 5}: Maximum Cardinality Bipartite Matching via Network Flow Reduction and Edmonds-Karp Max-Flow.
    \item \textbf{Assignment 6}: Multi-Criteria Airline Flight Itinerary Optimization using Dijkstra's Shortest Path Algorithm.
\end{enumerate}

\vspace{0.3cm}

% ==============================================================================
% ASSIGNMENT 1: SOFTWARE LICENSES
% ==============================================================================
\section*{Assignment 1: Optimal Cryptographic Software License Purchase Schedule}

\subsection*{1.1 Problem Statement and Mathematical Formulation}
\textbf{Problem}: A security consultancy needs to acquire licenses for $n$ distinct cryptographic software packages. Regulations restrict purchases to at most one license per month. At month $t = 0$, every license has a base price of $C_0 = \text{Rs. } 1000$. However, the price of license $j$ escalates at an exponential monthly rate $r_j > 1$. If license $j$ is bought $t$ months from now ($t \in \{0, 1, \dots, n-1\}$), its cost is:
\[
\text{Cost}(j, t) = C_0 \cdot r_j^t
\]
All growth rates $r_1, r_2, \dots, r_n$ are distinct positive real numbers strictly greater than 1. The objective is to determine a permutation $\pi$ of the $n$ licenses minimizing the total expenditure:
\[
\min_{\pi} \sum_{t=0}^{n-1} C_0 \cdot r_{\pi(t)}^t
\]

\textbf{Proof of Optimality (Greedy Choice via Exchange Argument)}:
\begin{itemize}
    \item \textbf{Greedy Rule}: Sort and purchase licenses in strictly non-increasing order of their monthly growth rates:
    \[
    r_{\pi(0)} > r_{\pi(1)} > \dots > r_{\pi(n-1)}
    \]
    \item \textbf{Exchange Argument}: Suppose an optimal schedule $\text{OPT}$ contains an inversion: there exist two consecutively purchased licenses $a$ at month $t$ and $b$ at month $t+1$ such that $r_a < r_b$.
    \item The joint cost contributed by these two licenses in $\text{OPT}$ is:
    \[
    C_{\text{OPT}} = C_0 \left( r_a^t + r_b^{t+1} \right)
    \]
    \item Invert the purchase order to yield schedule $\text{OPT}'$ (buying $b$ at month $t$ and $a$ at month $t+1$):
    \[
    C_{\text{OPT}'} = C_0 \left( r_b^t + r_a^{t+1} \right)
    \]
    \item Evaluating the cost difference:
    \[
    C_{\text{OPT}} - C_{\text{OPT}'} = C_0 \left[ (r_b^{t+1} - r_b^t) - (r_a^{t+1} - r_a^t) \right] = C_0 \left[ r_b^t(r_b - 1) - r_a^t(r_a - 1) \right]
    \]
    \item Since $r_b > r_a > 1$ and $t \ge 0$, it follows that $r_b^t \ge r_a^t > 0$ and $(r_b - 1) > (r_a - 1) > 0$.
    \item Therefore, $C_{\text{OPT}} - C_{\text{OPT}'} > 0 \implies C_{\text{OPT}'} < C_{\text{OPT}}$, strictly reducing the total cost.
    \item Because every inversion strictly increases the cost, the greedy sorted order must achieve the global minimum.
\end{itemize}

\subsection*{1.2 Python Implementation (\texttt{software\_licenses.py})}
\begin{lstlisting}[language=Python]
from dataclasses import dataclass
from typing import List, Tuple

@dataclass
class License:
    name: str
    rate: float
    base_cost: float = 1000.0

def optimal_license_schedule(
    licenses: List[License],
) -> Tuple[List[Tuple[int, License, float]], float]:
    sorted_licenses = sorted(licenses, key=lambda lic: lic.rate, reverse=True)
    schedule = []
    total_cost = 0.0

    for month, lic in enumerate(sorted_licenses):
        cost_at_t = lic.base_cost * (lic.rate ** month)
        schedule.append((month, lic, cost_at_t))
        total_cost += cost_at_t

    return schedule, total_cost
\end{lstlisting}

\subsection*{1.3 Theoretical Complexity Analysis}
\begin{itemize}
    \item \textbf{Sorting Time}: Sorting $n$ license rates takes $\mathcal{O}(n \log n)$ time.
    \item \textbf{Schedule Generation}: A single linear scan evaluates exponential costs in $\mathcal{O}(n)$ time.
    \item \textbf{Total Time Complexity}: $\mathcal{O}(n \log n)$.
    \item \textbf{Auxiliary Space Complexity}: $\mathcal{O}(n)$ to store the ordered schedule.
\end{itemize}

\subsection*{1.4 Execution Output \& Verification}
\begin{lstlisting}
Month  | License Name              | Rate (r_j)   | Cost (Rs.)     
-----------------------------------------------------------------
Month 0  | Post-Quantum Signer       | 1.80         | Rs. 1000.00     
Month 1  | RSA Key Manager           | 1.50         | Rs. 1500.00     
Month 2  | Zero-Knowledge Prover     | 1.35         | Rs. 1822.50     
Month 3  | AES Encryption Suite      | 1.20         | Rs. 1728.00     
Month 4  | Elliptic Curve Tool       | 1.10         | Rs. 1464.10     
-----------------------------------------------------------------
Minimum Total Purchase Cost: Rs. 7514.60
Verification Passed: Highest growth rates purchased earliest!
\end{lstlisting}

\newpage

% ==============================================================================
% ASSIGNMENT 2: EULERIAN CIRCUIT
% ==============================================================================
\section*{Assignment 2: Linear-Time Directed Eulerian Circuit Construction}

\subsection*{2.1 Problem Formulation and Existence Theorem}
\textbf{Definition}: An Eulerian circuit in a directed graph $G = (V, E)$ is a closed directed walk traversing every directed edge in $E$ exactly once.

\textbf{Euler's Theorem for Directed Graphs}:
A directed multigraph $G$ possesses an Eulerian circuit if and only if:
\begin{enumerate}
    \item \textbf{Degree Balance}: For every vertex $v \in V$, its in-degree equals its out-degree:
    \[
    d_{\text{in}}(v) = d_{\text{out}}(v) \quad \forall v \in V
    \]
    \item \textbf{Connectivity}: All vertices with non-zero degree ($d(v) > 0$) belong to a single connected component in the underlying undirected graph.
\end{enumerate}

\textbf{Hierholzer's Algorithm}:
\begin{enumerate}
    \item Choose an arbitrary starting vertex $s$ with $d_{\text{out}}(s) > 0$.
    \item Maintain a stack for the current exploratory path. Follow unused outgoing edges from the current vertex, pushing newly visited vertices onto the stack.
    \item When the vertex at the top of the stack has no remaining unused outgoing edges, pop it and append it to the circuit.
    \item Upon exhausting all edges, reverse the circuit to recover the traversal sequence in forward chronological order.
\end{enumerate}

\subsection*{2.2 Python Implementation (\texttt{eulerian\_circuit.py})}
\begin{lstlisting}[language=Python]
from collections import defaultdict
from typing import Dict, List, Optional, Tuple

class DirectedGraph:
    def __init__(self):
        self.adj = defaultdict(list)
        self.in_degree = defaultdict(int)
        self.out_degree = defaultdict(int)
        self.num_edges = 0

    def add_edge(self, u: int, v: int) -> None:
        self.adj[u].append(v)
        self.out_degree[u] += 1
        self.in_degree[v] += 1
        self.num_edges += 1

    def find_eulerian_circuit(self) -> Tuple[Optional[List[int]], str]:
        # 1. Degree balance and connectivity verification
        all_v = set(self.in_degree.keys()) | set(self.out_degree.keys())
        start = None
        for v in all_v:
            if self.in_degree[v] != self.out_degree[v]:
                return None, f"Vertex {v} unbalanced (in={self.in_degree[v]}, out={self.out_degree[v]})."
            if self.out_degree[v] > 0 and start is None:
                start = v

        if self.num_edges == 0: return [], "Trivially Eulerian."

        # 2. Hierholzer's Algorithm using backtrack stack
        edge_lists = {u: list(neighbors) for u, neighbors in self.adj.items()}
        circuit, stack = [], [start]

        while stack:
            curr = stack[-1]
            if edge_lists.get(curr):
                stack.append(edge_lists[curr].pop())
            else:
                circuit.append(stack.pop())

        circuit.reverse()
        if len(circuit) != self.num_edges + 1:
            return None, "Graph contains disconnected edge components."
        return circuit, "Eulerian circuit successfully constructed."
\end{lstlisting}

\subsection*{2.3 Complexity Analysis}
\begin{itemize}
    \item \textbf{Degree \& Connectivity Verification}: Incurs $\mathcal{O}(V + E)$ time.
    \item \textbf{Hierholzer's Traversal}: Each directed edge is pushed onto the call stack once and popped in $\mathcal{O}(1)$ time using adjacency lists.
    \item \textbf{Total Time Complexity}: $\mathcal{O}(V + E)$ (strictly optimal linear time).
    \item \textbf{Auxiliary Space Complexity}: $\mathcal{O}(V + E)$ to maintain edge tracking buffers and the traversal stack.
\end{itemize}

\subsection*{2.4 Execution Output \& Verification}
\begin{lstlisting}
Graph 1 Edges: [(0, 1), (1, 2), (2, 0), (0, 3), (3, 4), (4, 0)]
Status       : Eulerian circuit successfully constructed.
Euler Circuit: 0 -> 3 -> 4 -> 0 -> 1 -> 2 -> 0
Verification Passed: All edges visited exactly once in cycle!

Graph 2 Edges: [(0, 1), (1, 2), (2, 0), (2, 3)]
Status       : Vertex 2 unbalanced (in=1, out=2).
Verification Passed: Correctly identified non-Eulerian directed graph.
\end{lstlisting}

\newpage

% ==============================================================================
% ASSIGNMENT 3: PICKUP STICKS
% ==============================================================================
\section*{Assignment 3: 3D Pick-up Sticks Geometric Relation Routine \& Scheduling}

\subsection*{3.1 Algorithmic Formulation and 3D Line Intersection}
\textbf{Problem}: Given $N$ non-vertical sticks in 3D space, each defined by two endpoints $(x_1, y_1, z_1)$ and $(x_2, y_2, z_2)$:
\begin{enumerate}
    \item[(a)] Determine whether stick $a$ is \textbf{above}, \textbf{below}, or \textbf{unrelated} to stick $b$.
    \item[(b)] Determine whether all sticks can be picked up, and output a valid pickup sequence.
\end{enumerate}

\textbf{Geometric Routine}:
\begin{itemize}
    \item Parameterize stick $a$ and stick $b$ on the horizontal $xy$-plane:
    \[
    A(t) = A_1 + t(A_2 - A_1), \quad t \in [0, 1]
    \]
    \[
    B(u) = B_1 + u(B_2 - B_1), \quad u \in [0, 1]
    \]
    \item Solve the linear system for the 2D projection intersection $(t, u)$ via Cramer's rule:
    \[
    D = \Delta x_a \Delta y_b - \Delta y_a \Delta x_b
    \]
    \item If $|D| < 10^{-9}$ (segments parallel) or $t, u \notin [0, 1]$, the sticks do not intersect in the top-down projection $\implies$ \textbf{unrelated}.
    \item If $t, u \in [0, 1]$, compute the height of both sticks at the intersection point:
    \[
    z_a = z_{A1} + t(z_{A2} - z_{A1}), \qquad z_b = z_{B1} + u(z_{B2} - z_{B1})
    \]
    \item If $z_a > z_b$, stick $a$ rests on top of stick $b$ ($\mathbf{a \text{ is above } b}$).
    \item If $z_a < z_b$, stick $a$ rests below stick $b$ ($\mathbf{a \text{ is below } b}$).
\end{itemize}

\textbf{Topological Scheduling}:
A stick can only be picked up when no other stick lies on top of it. If stick $a$ is above stick $b$, $a$ must be removed before $b$, creating a directed precedence constraint $a \to b$. We run Kahn's algorithm: sticks with in-degree 0 (no sticks on top) are repeatedly picked up. If all $N$ sticks are eliminated, the sequence is valid; otherwise, a cycle exists (deadlock).

\subsection*{3.2 Python Implementation (\texttt{pickup\_sticks.py})}
\begin{lstlisting}[language=Python]
from collections import defaultdict, deque
from dataclasses import dataclass
from typing import Dict, List, Optional, Tuple

@dataclass
class Stick:
    id: int
    p1: Tuple[float, float, float]
    p2: Tuple[float, float, float]

def check_stick_relation(a: Stick, b: Stick, eps: float = 1e-9) -> str:
    dx_a, dy_a = a.p2[0] - a.p1[0], a.p2[1] - a.p1[1]
    dx_b, dy_b = b.p2[0] - b.p1[0], b.p2[1] - b.p1[1]
    denom = dx_a * dy_b - dy_a * dx_b

    if abs(denom) < eps: return "unrelated"

    dx_ab, dy_ab = b.p1[0] - a.p1[0], b.p1[1] - a.p1[1]
    t = (dx_ab * dy_b - dy_ab * dx_b) / denom
    u = (dx_ab * dy_a - dy_ab * dx_a) / denom

    if -eps <= t <= 1.0 + eps and -eps <= u <= 1.0 + eps:
        t_c, u_c = max(0.0, min(1.0, t)), max(0.0, min(1.0, u))
        z_a = a.p1[2] + t_c * (a.p2[2] - a.p1[2])
        z_b = b.p1[2] + u_c * (b.p2[2] - b.p1[2])
        if z_a > z_b + eps: return "above"
        elif z_a < z_b - eps: return "below"
    return "unrelated"

def find_pickup_sequence(sticks: List[Stick]) -> Tuple[bool, Optional[List[int]], str]:
    n = len(sticks)
    adj = defaultdict(list)
    in_degree = {s.id: 0 for s in sticks}

    for i in range(n):
        for j in range(i + 1, n):
            rel = check_stick_relation(sticks[i], sticks[j])
            if rel == "above":
                adj[sticks[i].id].append(sticks[j].id)
                in_degree[sticks[j].id] += 1
            elif rel == "below":
                adj[sticks[j].id].append(sticks[i].id)
                in_degree[sticks[i].id] += 1

    queue = deque([sid for sid, deg in in_degree.items() if deg == 0])
    order = []
    while queue:
        curr = queue.popleft()
        order.append(curr)
        for nbr in adj[curr]:
            in_degree[nbr] -= 1
            if in_degree[nbr] == 0: queue.append(nbr)

    if len(order) == n: return True, order, "Valid sequence found."
    return False, None, "Cycle detected: Deadlock prevents full pickup."
\end{lstlisting}

\subsection*{3.3 Complexity Analysis}
\begin{itemize}
    \item \textbf{Pairwise Intersection Routine}: Evaluating all pairs of $N$ sticks requires $\binom{N}{2} = \mathcal{O}(N^2)$ checks.
    \item \textbf{Topological Sort}: Kahn's algorithm operates on $V = N$ vertices and $E \le N^2$ edges in $\mathcal{O}(V + E) = \mathcal{O}(N^2)$ time.
    \item \textbf{Total Time Complexity}: $\mathcal{O}(N^2)$.
    \item \textbf{Space Complexity}: $\mathcal{O}(N^2)$ to store the dependency graph.
\end{itemize}

\subsection*{3.4 Execution Output \& Verification}
\begin{lstlisting}
--- (a) Pairwise Relation Routine Demonstration ---
Relation(Stick 1, Stick 2): above (Expected: above)
Relation(Stick 2, Stick 3): above (Expected: above)
Relation(Stick 1, Stick 4): unrelated (Expected: unrelated)

--- (b) Pickup Sequence Determination ---
Status        : Valid pickup sequence found.
Pickup Order  : [1, 4, 2, 3]
Verification Passed: Valid topological stick pickup sequence generated!
\end{lstlisting}

\newpage

% ==============================================================================
% ASSIGNMENT 4: TWO-COLORABILITY
% ==============================================================================
\section*{Assignment 4: Linear-Time Two-Colorability Testing \& Applications}

\subsection*{4.1 Theoretical Formulation and Bipartite Relationship}
\textbf{Definition}: A graph $G = (V, E)$ is $k$-colorable if each vertex can be assigned one of $k$ colors such that no edge connects two vertices of identical color.

\textbf{Relationship to Bipartite Graphs}:
\begin{itemize}
    \item \textbf{Equivalence Theorem}: A graph is \textbf{2-colorable if and only if it is bipartite}.
    \item \textbf{Odd Cycle Characterization}: A graph is 2-colorable / bipartite if and only if it contains \textbf{no odd-length cycles}.
    \item If a graph is 2-colorable with colors $\{0, 1\}$, partitioning vertices by color yields two disjoint sets $L = \{v \mid c(v) = 0\}$ and $R = \{v \mid c(v) = 1\}$. Every edge has one endpoint in $L$ and the other in $R$, satisfying the definition of a bipartite graph.
\end{itemize}

\textbf{Two Practical Applications}:
\begin{enumerate}
    \item \textbf{Maximum Matching and Assignment Problems}: Matching medical residents to hospitals, organ donors to recipients, or distributed compute jobs to worker nodes requires modeling conflict-free pairings over bipartite graphs.
    \item \textbf{Conflict-Free Frequency and Register Allocation}: Determining whether wireless transmitters, exam timetables, or compiler variable registers can operate without interference across two alternating channels or time slots.
\end{enumerate}

\subsection*{4.2 Python Implementation (\texttt{two\_colorability.py})}
\begin{lstlisting}[language=Python]
from collections import defaultdict, deque
from typing import Dict, List, Optional, Tuple

def test_two_colorability(
    adj: Dict[int, List[int]], vertices: List[int]
) -> Tuple[bool, Optional[Dict[int, int]], Optional[Tuple[List[int], List[int]]], str]:
    colors: Dict[int, int] = {v: -1 for v in vertices}

    for start in vertices:
        if colors[start] != -1: continue
        colors[start] = 0
        queue = deque([start])

        while queue:
            u = queue.popleft()
            for v in adj[u]:
                if colors[v] == -1:
                    colors[v] = 1 - colors[u]
                    queue.append(v)
                elif colors[v] == colors[u]:
                    return False, None, None, f"Odd cycle: ({u}, {v}) share color {colors[u]}."

    set_L = [v for v, c in colors.items() if c == 0]
    set_R = [v for v, c in colors.items() if c == 1]
    return True, colors, (set_L, set_R), "Graph is 2-colorable (bipartite)."
\end{lstlisting}

\subsection*{4.3 Complexity Analysis}
\begin{itemize}
    \item \textbf{Time Complexity}: $\mathcal{O}(V + E)$ using Breadth-First Search (BFS) over adjacency lists.
    \item \textbf{Space Complexity}: $\mathcal{O}(V)$ for the color map and BFS queue.
\end{itemize}

\subsection*{4.4 Execution Output \& Verification}
\begin{lstlisting}
Graph 1 (4-Cycle): Graph is 2-colorable (bipartite).
 - Color Assignments: {1: 0, 2: 1, 3: 0, 4: 1}
 - Bipartite Sets   : Set L = [1, 3], Set R = [2, 4]
Verification Passed: Correctly identified 2-colorable graph!

Graph 2 (Triangle / Odd Cycle): Odd cycle: (2, 3) share color 1.
Verification Passed: Correctly rejected odd-cycle graph!
\end{lstlisting}

\newpage

% ==============================================================================
% ASSIGNMENT 5: BIPARTITE MATCHING VIA FLOW
% ==============================================================================
\section*{Assignment 5: Maximum Bipartite Matching via Max-Flow Reduction}

\subsection*{5.1 Network Flow Reduction Principle}
\textbf{Problem}: Given a bipartite graph $G = (L \cup R, E)$ where edges only connect vertices between independent sets $L$ and $R$, find the maximum-cardinality matching.

\textbf{Construction of Flow Network $G'$}:
\begin{enumerate}
    \item Introduce a super-source $s$ and a super-sink $t$.
    \item Construct directed edges $(s, u)$ with capacity $c(s, u) = 1$ for all $u \in L$.
    \item Construct directed edges $(u, v)$ with capacity $c(u, v) = 1$ for all original edges $(u, v) \in E$ with $u \in L, v \in R$.
    \item Construct directed edges $(v, t)$ with capacity $c(v, t) = 1$ for all $v \in R$.
\end{enumerate}

\textbf{Integrality Theorem and Equivalence}:
Because all capacities are integral (equal to 1), the Integrality Theorem guarantees the existence of a maximum flow where every edge carries an integer flow $f(e) \in \{0, 1\}$.
\begin{itemize}
    \item At most 1 unit of flow leaves $s$ into each $u \in L$, guaranteeing each left vertex is matched at most once.
    \item At most 1 unit of flow enters $t$ from each $v \in R$, guaranteeing each right vertex is matched at most once.
    \item By flow conservation, the set of edges $(u, v) \in E$ with $f(u, v) = 1$ constitutes a valid matching.
    \item Therefore, $\text{Value}(\text{Max-Flow}) = \text{Cardinality}(\text{Max-Matching})$.
\end{itemize}

\subsection*{5.2 Python Implementation (\texttt{bipartite\_matching.py})}
\begin{lstlisting}[language=Python]
from collections import defaultdict, deque
from typing import Dict, List, Set, Tuple

class FlowNetwork:
    def __init__(self):
        self.adj = defaultdict(list)
        self.capacity = defaultdict(int)
        self.flow = defaultdict(int)

    def add_edge(self, u: str, v: str, cap: int):
        self.adj[u].append(v)
        self.adj[v].append(u)
        self.capacity[(u, v)] = cap
        self.capacity[(v, u)] = 0

    def edmonds_karp_max_flow(self, source: str, sink: str) -> int:
        max_flow = 0
        parent = {}

        while True:
            visited = {source}
            queue = deque([source])
            found = False
            while queue and not found:
                curr = queue.popleft()
                for nbr in self.adj[curr]:
                    residual = self.capacity[(curr, nbr)] - self.flow[(curr, nbr)]
                    if nbr not in visited and residual > 0:
                        visited.add(nbr)
                        parent[nbr] = curr
                        if nbr == sink: found = True; break
                        queue.append(nbr)
            if not found: break

            path_flow = min(self.capacity[(parent[c], c)] - self.flow[(parent[c], c)] 
                            for c in parent if c == sink or c in parent.values())
            curr = sink
            while curr != source:
                p = parent[curr]
                self.flow[(p, curr)] += 1
                self.flow[(curr, p)] -= 1
                curr = p
            max_flow += 1
            parent.clear()
        return max_flow
\end{lstlisting}

\subsection*{5.3 Complexity Analysis}
\begin{itemize}
    \item \textbf{Edmonds-Karp Complexity}: In general networks, Edmonds-Karp runs in $\mathcal{O}(V \cdot E^2)$.
    \item \textbf{Unit Capacity Bipartite Bound}: Because unit edge capacities bound the maximum flow by $\min(|L|, |R|) \le V$, at most $V$ augmenting path augmentations occur. Each BFS takes $\mathcal{O}(E)$ time.
    \item \textbf{Total Time Complexity}: $\mathcal{O}(V \cdot E)$.
    \item \textbf{Space Complexity}: $\mathcal{O}(V + E)$ for residual graph storage.
\end{itemize}

\subsection*{5.4 Execution Output \& Verification}
\begin{lstlisting}
Set L (Applicants) : ['Alice', 'Bob', 'Charlie', 'Diana']
Set R (Job Roles)  : ['Software Engineer', 'Data Scientist', 'Cybersecurity Analyst', 'DevOps Engineer']
Available Edges    : 7 candidate qualifications

Maximum Matching Size (Max-Flow): 4
Optimal Assigned Pairs:
   * Alice      <---> Software Engineer
   * Bob        <---> Cybersecurity Analyst
   * Charlie    <---> Data Scientist
   * Diana      <---> DevOps Engineer
Verification Passed: Valid maximum cardinality bipartite matching obtained!
\end{lstlisting}

\newpage

% ==============================================================================
% ASSIGNMENT 6: FLIGHT SCHEDULING
% ==============================================================================
\section*{Assignment 6: Airline Flight Route Optimization (Dijkstra's Algorithm)}

\subsection*{6.1 Problem Statement and Algorithmic Design}
\textbf{Problem}: A traveler wishes to journey from origin city $A$ to destination city $B$. Flight networks are modeled as a directed graph where vertices represent metropolitan airports, and directed edges represent scheduled flights endowed with dual cost metrics:
\begin{enumerate}
    \item \textbf{Monetary Cost}: Fare in Rupees ($C_{uv} > 0$).
    \item \textbf{Travel Time}: Flight flight duration in hours ($T_{uv} > 0$).
\end{enumerate}
The algorithm must determine the shortest travel itinerary minimizing either total cost or total flight duration using Dijkstra's single-source shortest path algorithm.

\subsection*{6.2 Python Implementation (\texttt{flight\_scheduling.py})}
\begin{lstlisting}[language=Python]
import heapq
from collections import defaultdict
from dataclasses import dataclass
from typing import Dict, List, Optional, Tuple

@dataclass
class Flight:
    origin: str
    destination: str
    airline: str
    flight_num: str
    duration_hours: float
    cost_inr: float

class FlightNetwork:
    def __init__(self):
        self.adj = defaultdict(list)
        self.cities = set()

    def add_flight(self, orig, dest, airline, fnum, dur, cost):
        f = Flight(orig, dest, airline, fnum, dur, cost)
        self.adj[orig].append(f)
        self.cities.update([orig, dest])

    def find_shortest_itinerary(
        self, source: str, target: str, criterion: str = "cost"
    ) -> Tuple[Optional[float], Optional[List[Flight]], str]:
        pq = [(0.0, source)]
        dist = {source: 0.0}
        prev_flight = {}

        while pq:
            d_curr, curr = heapq.heappop(pq)
            if d_curr > dist.get(curr, float("inf")): continue
            if curr == target: break

            for flight in self.adj[curr]:
                w = flight.cost_inr if criterion == "cost" else flight.duration_hours
                if d_curr + w < dist.get(flight.destination, float("inf")):
                    dist[flight.destination] = d_curr + w
                    prev_flight[flight.destination] = (curr, flight)
                    heapq.heappush(pq, (d_curr + w, flight.destination))

        if target not in dist: return None, None, "No viable route."

        itinerary, curr = [], target
        while curr != source:
            prev_city, fl = prev_flight[curr]
            itinerary.append(fl)
            curr = prev_city
        itinerary.reverse()
        return dist[target], itinerary, f"Optimal route found for {criterion}."
\end{lstlisting}

\subsection*{6.3 Complexity Analysis}
\begin{itemize}
    \item \textbf{Time Complexity}: Using a min-priority queue (min-heap), Dijkstra's algorithm processes each vertex and relaxes each directed flight edge in $\mathcal{O}((V + E) \log V)$ time.
    \item \textbf{Auxiliary Space Complexity}: $\mathcal{O}(V + E)$ to maintain graph adjacency and predecessor pointers.
\end{itemize}

\subsection*{6.4 Execution Output \& Verification}
\begin{lstlisting}
Journey Planning: Jammu ---> Bengaluru

--- 1. Optimization Criterion: MINIMUM COST (INR) ---
Status: Optimal route found optimizing for cost.
Total Cost: Rs. 10700.00
Flight Route:
  Leg 1: Jammu -> Delhi via Air India (AI-822) | Duration: 1.5h | Cost: Rs. 4500.00
  Leg 2: Delhi -> Bengaluru via IndiGo (6E-501) | Duration: 2.8h | Cost: Rs. 6200.00
Total Flight Time: 4.3 hours

--- 2. Optimization Criterion: MINIMUM TRAVEL TIME (Hours) ---
Status: Optimal route found optimizing for time.
Total Duration: 4.3 hours
Flight Route:
  Leg 1: Jammu -> Delhi via Air India (AI-822) | Duration: 1.5h | Cost: Rs. 4500.00
  Leg 2: Delhi -> Bengaluru via IndiGo (6E-501) | Duration: 2.8h | Cost: Rs. 6200.00
Total Monetary Cost: Rs. 10700.00
Verification Passed: Dijkstra shortest-path routes successfully computed!
\end{lstlisting}

\newpage

% ==============================================================================
% SUMMARY AND CONCLUSION
% ==============================================================================
\section*{Laboratory Summary \& Comprehensive Complexity Bounds}

Table \ref{tab:lab9_summary} aggregates the core algorithmic classifications, optimal paradigms, and theoretical complexity bounds established across all six assignments.

\begin{table}[H]
\centering
\caption{Comprehensive Algorithmic Bounds for Lab Assignment 9.}
\label{tab:lab9_summary}
\vspace{0.2cm}
\resizebox{\textwidth}{!}{
\begin{tabular}{lllll}
\toprule
\textbf{Assignment} & \textbf{Problem Domain} & \textbf{Algorithmic Paradigm} & \textbf{Time Complexity} & \textbf{Space Complexity} \\
\midrule
\textbf{Assignment 1} & Software License Scheduling  & Greedy Strategy (Exchange Argument)  & $\mathcal{O}(n \log n)$            & $\mathcal{O}(n)$ \\
\textbf{Assignment 2} & Directed Eulerian Circuit    & Hierholzer's Backtrack Traversal     & $\mathcal{O}(V + E)$               & $\mathcal{O}(V + E)$ \\
\textbf{Assignment 3} & 3D Pick-up Sticks Routine    & Line Intersection \& Kahn's DAG Sort & $\mathcal{O}(N^2)$                 & $\mathcal{O}(N^2)$ \\
\textbf{Assignment 4} & Two-Colorability Testing     & Breadth-First Search (BFS)           & $\mathcal{O}(V + E)$               & $\mathcal{O}(V)$ \\
\textbf{Assignment 5} & Maximum Bipartite Matching   & Network Flow Reduction (Edmonds-Karp)& $\mathcal{O}(V \cdot E)$           & $\mathcal{O}(V + E)$ \\
\textbf{Assignment 6} & Flight Route Optimization    & Dijkstra's Shortest Path             & $\mathcal{O}((V + E) \log V)$      & $\mathcal{O}(V + E)$ \\
\bottomrule
\end{tabular}
}
\end{table}

\subsection*{Conclusion}
The algorithmic investigations conducted across Lab Assignment 9 illustrate several fundamental principles:
\begin{enumerate}
    \item \textbf{Power of Exchange Arguments}: Greedy rate sorting in software license purchases provably minimizes exponential cost divergence through local pairwise exchange.
    \item \textbf{Optimal Linear-Time Graph Traversals}: Both Eulerian circuit detection (Hierholzer's) and 2-colorability testing (alternating BFS) run in optimal $\mathcal{O}(V + E)$ time, bypassing brute-force path enumeration.
    \item \textbf{Bridging Geometry and Graph Scheduling}: Resolving 3D spatial overlaps into directed dependency arcs allows topological sorting to determine valid sequential actions or detect deadlock conditions.
    \item \textbf{Reductions to Network Flow}: Complex structural matching problems (bipartite matching) can be seamlessly modeled and resolved through polynomial-time max-flow cuts, grounded by the Integrality Theorem.
\end{enumerate}

\end{document}
